% ============================================================
%  PLANTILLA DE EXAMEN DE PRUEBA — syntheca
%  Clase ARTICLE (documento independiente, no un capítulo del
%  libro de la materia — un examen de prueba se guarda aparte,
%  en examenes-prueba/<materia>/).
%
%  Reutiliza la misma paleta y tipografía que la plantilla del
%  libro (plantilla-syntheca-preamble.tex) para que la identidad
%  visual sea consistente en todo lo que genera el sistema.
%
%  Diferencia clave de diseño respecto al libro: el entorno
%  \pregunta NO muestra metadata (familia, tipo 6.a/6.b, etc.)
%  — un examen de prueba tiene que verse como un examen real,
%  no como un documento de estudio con etiquetas de trazabilidad.
%  Esa trazabilidad vive en el nombre del archivo y en
%  mapa-estudio.json, no en el PDF que el estudiante rinde.
% ============================================================
\usepackage[utf8]{inputenc}
\usepackage[T1]{fontenc}
\usepackage{geometry}
\usepackage{amsmath}
\usepackage{amssymb}
\usepackage{enumitem}
\usepackage{tabularx}
\usepackage{booktabs}
\renewcommand{\familydefault}{\sfdefault}
\usepackage{titlesec}
\usepackage{fancyhdr}
\usepackage{ragged2e}
\usepackage[table]{xcolor}
\usepackage{array}

\geometry{
    top=2.2cm,
    bottom=2.2cm,
    left=2.2cm,
    right=2.2cm
}

\definecolor{textgray}{RGB}{60, 60, 60}
\definecolor{rulegray}{RGB}{150, 150, 150}
\color{textgray}

% ---------- Metadatos editables por examen ----------
\newcommand{\miNombre}{Nombre Apellido}
\newcommand{\materiaActual}{Materia}
\newcommand{\setMateria}[1]{\renewcommand{\materiaActual}{#1}}
\newcommand{\nombreExamen}{Parcial de práctica}
\newcommand{\setNombreExamen}[1]{\renewcommand{\nombreExamen}{#1}}
\newcommand{\duracionExamen}{}
\newif\ifhayduracion
\hayduracionfalse
\newcommand{\setDuracion}[1]{\renewcommand{\duracionExamen}{#1}\hayduraciontrue}

\pagestyle{fancy}
\fancyhf{}
\renewcommand{\headrulewidth}{0pt}
\renewcommand{\footrulewidth}{0pt}
\fancyfoot[C]{\small\color{textgray}\thepage}
\fancyfoot[R]{\footnotesize\color{textgray}\miNombre}
\fancyfoot[L]{\footnotesize\color{textgray}\materiaActual}

\renewcommand{\tablename}{Tabla}
\renewcommand{\figurename}{Figura}

% ---------- Encabezado del examen (no es \chapter — es un bloque
% de portada simple, el examen es un documento corto) ----------
\newcommand{\encabezadoExamen}{%
  \begin{center}
    {\fontsize{16}{19}\selectfont\bfseries\color{textgray} \MakeUppercase{\materiaActual}\par}
    \vspace{2pt}
    {\large\color{textgray} \nombreExamen \par}
    \vspace{4pt}
    {\color{rulegray}\rule{0.6\linewidth}{0.4pt}\par}
    \vspace{6pt}
    \ifhayduracion
      {\footnotesize\itshape\color{rulegray} Duración sugerida: \duracionExamen \par}
    \fi
    \vspace{10pt}
  \end{center}
}

% ---------- SECTION — se usa poco (a lo sumo para separar partes
% del examen, ej. "Parte A" / "Parte B"), mismo estilo que el libro ----------
\titleformat{\section}
  {\centering\color{textgray}\fontsize{12}{14}\selectfont\bfseries}
  {}
  {0pt}
  {}
  [\vspace{2pt}\color{rulegray}\titlerule\vspace{4pt}]
\titlespacing*{\section}{0pt}{14pt}{6pt}

% ============================================================
%  ENTORNO \pregunta — la pieza central de esta plantilla.
%  Numerado automático, SIN mostrar familia/tipo (a diferencia
%  de \ejercicio del libro) — argumento opcional solo para puntaje,
%  que es información legítima de un examen real.
%  Uso: \begin{pregunta}[25 puntos] ... \end{pregunta}
% ============================================================
\newcounter{preguntaexamen}
\newenvironment{pregunta}[1][]{%
  \refstepcounter{preguntaexamen}%
  \par\vspace{10pt}\noindent%
  \textbf{\color{textgray}Ejercicio \arabic{preguntaexamen}.}%
  \ifx&#1&\else\ {\footnotesize\itshape\color{rulegray}(#1)}\fi%
  \par\vspace{4pt}%
}{%
  \vspace{4pt}
}

% ============================================================
%  ENTORNO \resolucion — SOLO se usa en el documento de
%  resolución (nunca en el de enunciado, para no arruinar el
%  simulacro). Mismo criterio minimalista: una línea fina, sin caja.
% ============================================================
\newenvironment{resolucion}{%
  \par\vspace{4pt}\noindent%
  {\footnotesize\bfseries\color{rulegray}RESOLUCIÓN}\par%
  \color{rulegray}\hrule\color{textgray}\vspace{4pt}%
}{%
  \vspace{2pt}
}

\usepackage{ifthen}
